\chapter{SCP Regression Vectors and Integration Campaigns}
\label{app:scp03-vectors}

This appendix documents the SCP regression evidence used by the Oxide SE test
suite. It complements Chapter~3, Section~\ref{sec:protocol-scp}, Chapter~8,
Section~\ref{sec:impl-scp}, and Appendix~\ref{app:scp-data-objects}. The
validation campaign covers SCP03 S8, SCP03 S16, identity/no-SCP mode, and the
SCP11a/SCP11b/SCP11c establishment and secure-messaging paths. In addition to
the independently reproduced SCP03 files, it includes one byte-for-byte
SCP11a reference scenario published by Samsung.

\section{SCP03 Vector Files}

The deterministic SCP03 vector files validate the S8 and S16 profiles.

The files are:
\begin{itemize}
\item \texttt{xtask/tests/fixtures/scp03\_s8\_independently\_reproduced.txt};
\item \texttt{xtask/tests/fixtures/scp03\_s16\_independently\_reproduced.txt}.
\end{itemize}

They are documented in the repository as:

\begin{quote}
Reference vectors independently reproduced from OpenSCDP and GlobalPlatformPro
behavior, cross-checked against the public SCP03 specification.
\end{quote}

These are therefore \emph{not} official GlobalPlatform certification vectors.
They are independent regression vectors used to make sure that the Oxide SE
SCP03 implementation is not only self-consistent with its own host-side oracle.

\section{Samsung SCP11a Reference Scenario}

The first independent SCP11 fixture imported by Oxide SE is:

\begin{verbatim}
xtask/tests/fixtures/scp11a_p256_aes128_s16_samsung.txt
\end{verbatim}

It reproduces Samsung's SCP11a scenario using P-256, AES-128 and S16,
published in the Apache-2.0-licensed OpenSCP-Java project
\cite{samsung-openscp-java,samsung-openscp-blog}. The imported values come
from repository commit
\path{b9876fc36a5b18fb90ce03d0894f39edb08a905b}, specifically from:

\begin{itemize}
\item \path{Scp11TestData.java}, for the deterministic OCE static and
      ephemeral keys;
\item \path{Scp11Nist256TestData.java}, for the SD certificate material and
      certificate-management APDUs;
\item \path{SmartCardScp11aP256Aes128S16ModeEmulation.java}, for
      \texttt{MUTUAL AUTHENTICATE} and the first protected
      \texttt{GET STATUS} exchange.
\end{itemize}

The hexadecimal values are retained byte-for-byte. They are third-party
regression evidence, not GlobalPlatform certification vectors. Samsung
explicitly reports that public SCP11 vectors were not available and that its
team calculated these scenarios for OpenSCP-Java.

This scenario uses SCP identifier \texttt{11 01}: SCP11a without the optional
Host/Card identity binding. Its \texttt{SharedInfo} is therefore exactly
\texttt{3C 88 10}. The test validates:

\begin{itemize}
\item the normative top-level layout
      \texttt{A6 \ldots{} 5F49 \ldots{}} of the authentication data;
\item deterministic \texttt{ShSee} and \texttt{ShSss} ECDH results;
\item X9.63/SHA-256 key derivation over \texttt{ShSee || ShSss};
\item the Amendment-F assignment order: receipt key, \texttt{S-ENC},
      \texttt{S-MAC}, \texttt{S-RMAC}, then \texttt{S-DEK};
\item the receipt over the exact command \texttt{A6}/\texttt{5F49} objects
      followed by the card \texttt{5F49} object;
\item the first S16 protected \texttt{GET STATUS} C-APDU and R-APDU published
      by Samsung, including R-MAC verification and response decryption.
\end{itemize}

The corresponding host-based regression command is:

\begin{verbatim}
cargo test -p xtask --test scp11_vectors --offline
\end{verbatim}

The protected Samsung C-APDU is reproduced by the fixture test to prove the
derived keys and MAC chain. It is not injected unchanged into the QEMU
campaign: oXiDe SE kernel integration continues to exercise the common
Amendment~D framing selected for its SCP03/SCP11 APDU layer. The independent fixture
and the QEMU campaigns consequently validate complementary boundaries.

\section{Why Two Profiles Are Tested}

Oxide SE currently validates two explicit SCP03 profiles in its kernel-native
Security Domain path:
\begin{itemize}
\item \texttt{S8};
\item \texttt{S16}.
\end{itemize}

The S8 profile uses:
\begin{itemize}
\item 8-byte host and card challenges;
\item 8-byte authentication cryptograms;
\item 8-byte truncated C-MAC and R-MAC values.
\end{itemize}

The S16 profile uses:
\begin{itemize}
\item 16-byte host and card challenges;
\item 16-byte authentication cryptograms;
\item 16-byte non-truncated C-MAC and R-MAC values.
\end{itemize}

Testing both profiles is important because S16 is not only a larger version of
S8. It forces the APDU secure-channel layer to stop assuming that a protected
command always ends with an 8-byte MAC. In oXiDe SE, the APDU layer asks the
active Security Domain for the current MAC length before separating the raw
command payload from its trailing C-MAC. The S16 tests are therefore also
integration tests for that profile-dependent boundary.

\section{Vector File Structure}

Each vector file is a small line-oriented key/value document. Comments describe
the profile and provenance. The remaining fields are hexadecimal byte strings.

The common fields are:
\begin{itemize}
\item \texttt{profile}: either \texttt{S8} or \texttt{S16};
\item \texttt{static\_enc\_key}: the AES-128 static ENC key used for session
      derivation;
\item \texttt{static\_mac\_key}: the AES-128 static MAC key used for session
      derivation and cryptograms;
\item \texttt{host\_challenge}: the host challenge sent in
      \texttt{INITIALIZE UPDATE};
\item \texttt{card\_challenge}: the deterministic development card challenge;
\item \texttt{context}: the SCP03 derivation context, formed from host and
      card challenges.
\end{itemize}

The authentication and session-derivation fields are:
\begin{itemize}
\item \texttt{card\_cryptogram};
\item \texttt{host\_cryptogram};
\item \texttt{session\_enc\_key};
\item \texttt{session\_mac\_key};
\item \texttt{session\_rmac\_key}.
\end{itemize}

The secure-messaging fields exercise command and response protection:
\begin{itemize}
\item \texttt{protected\_get\_data\_header};
\item \texttt{protected\_get\_data\_cmac};
\item \texttt{protected\_get\_data\_response};
\item \texttt{protected\_get\_data\_rmac};
\item \texttt{response\_enc\_iv\_1};
\item \texttt{encrypted\_lifecycle\_response};
\item \texttt{encrypted\_lifecycle\_rmac}.
\end{itemize}

The command-encryption fields exercise \texttt{C-ENC}:
\begin{itemize}
\item \texttt{store\_plaintext};
\item \texttt{store\_enc\_iv\_1};
\item \texttt{store\_encrypted\_data};
\item \texttt{store\_cmac}.
\end{itemize}

Finally, the protected installation fields exercise a GlobalPlatform
\texttt{INSTALL [for install]} payload:
\begin{itemize}
\item \texttt{install\_plain\_data};
\item \texttt{install\_enc\_iv\_2};
\item \texttt{install\_encrypted\_data};
\item \texttt{install\_cmac}.
\end{itemize}

\section{How The Tests Use The Vectors}

The pure SCP03 vector test is:

\begin{verbatim}
cargo test -p xtask --test scp03_vectors --offline
\end{verbatim}

This test parses both files and recomputes the expected values with the
host-side crypto implementation used by the test oracle. It checks:
\begin{itemize}
\item KDF output for card and host cryptograms;
\item session ENC, MAC and RMAC keys;
\item protected \texttt{GET DATA} C-MAC;
\item protected response R-MAC;
\item response encryption and its R-MAC;
\item encrypted \texttt{STORE DATA} payload and C-MAC;
\item encrypted \texttt{INSTALL [for install]} payload and C-MAC.
\end{itemize}

The dedicated SCP03 QEMU integration test is:

\begin{verbatim}
cargo run test gp_scp03 --elf s8 mps2-an385
\end{verbatim}

It boots the kernel-native SCP03 profiles:
\begin{itemize}
\item one with \texttt{KernelSecurityDomain} plus
      \texttt{OXIDE\_SE\_SCP03\_PROFILE=S8};
\item one with \texttt{KernelSecurityDomain} plus
      \texttt{OXIDE\_SE\_SCP03\_PROFILE=S16}.
\end{itemize}

For the two kernel SCP03 profiles, the test performs the current management
regression flow:
\begin{enumerate}
\item clear \texttt{GET DATA};
\item \texttt{INITIALIZE UPDATE};
\item bad \texttt{EXTERNAL AUTHENTICATE} rejection;
\item successful \texttt{EXTERNAL AUTHENTICATE};
\item bad C-MAC rejection;
\item protected \texttt{GET DATA} with R-MAC;
\item clear \texttt{INSTALL [for install]} rejection;
\item encrypted secure-channel establishment;
\item protected \texttt{GET DATA} with R-ENC and R-MAC;
\item encrypted \texttt{STORE DATA};
\item encrypted \texttt{PUT KEY} and rotated-keyset authentication;
\item encrypted \texttt{INSTALL [for install]};
\item secondary kernel Security Domain instantiation and non-escalation checks;
\item selection and execution of \texttt{minimal\_valid\_test}.
\end{enumerate}

The identity/no-SCP path is validated separately through:

\begin{verbatim}
cargo run test gp_noscp --elf mps2-an385
\end{verbatim}

Oxide SE also maintains a dedicated Rustlet-backed Security Domain campaign.
That complementary path validates the same SCP03 concepts through a selected
Security Domain instance reached via \texttt{RustletSecurityDomainProxy}. This
matters because it demonstrates that the same secure-channel framework can be
driven either by a kernel-native backend or by a user-land Rustlet authority.
That campaign also validates that SCP03 keysets stored through
\texttt{PUT KEY} are scoped per Security Domain instance: a rotated keyset
installed for one Rustlet-backed Security Domain instance is not visible from a
second instance until that second instance receives its own \texttt{PUT KEY}.

The SCP11 establishment and secure-messaging smoke paths are validated through:

\begin{verbatim}
cargo run test gp_scp11a --elf mps2-an385
cargo run test gp_scp11b --elf mps2-an385
cargo run test gp_scp11c --elf mps2-an385
\end{verbatim}

Those tests exercise \texttt{PERFORM SECURITY OPERATION} where the profile
requires it. They also validate \texttt{MUTUAL AUTHENTICATE} and
\texttt{INTERNAL AUTHENTICATE}, the TLV-coded \texttt{A6} request, the
\texttt{5F49}/\texttt{86} response, receipt verification, and protected
payload-and-MAC secure messaging through the generic secure-channel layer. They prove
the internally selected development profiles and their dispatch paths; they do
not by themselves prove that every establishment parameter and cryptographic
role matches the complete GlobalPlatform profile.

The vector files and the QEMU tests serve different purposes. The vector files
validate deterministic SCP03 transformations and the independently published
Samsung SCP11a establishment scenario. The QEMU tests
validate that SCP establishment and secure messaging are correctly connected to
APDU parsing, Security-Domain policy, registry mutation, and Rustlet execution.

\section{Provenance and Limits}

OpenSCDP and GlobalPlatformPro are used as independent behavioral references
for reproducing the vectors. They are valuable because they exercise the same
GlobalPlatform/SCP03 concepts outside the Oxide SE implementation. The vectors
were also cross-checked against the public SCP03 specification.

For SCP11a, Samsung OpenSCP-Java is the independent implementation and source
of the fixed S16 scenario. Its repository provides the certificates, static
and ephemeral keys, C-APDUs and R-APDUs needed to reproduce that scenario
without executing Oxide SE's own host oracle. No equivalent Samsung SCP11b
fixture is present in the referenced test suite.

This still leaves an explicit certification gap:
\begin{itemize}
\item the files are not official GlobalPlatform certification vectors;
\item development bootstrap material still exists for the current kernel
      profiles and must be replaced by a production provisioning policy;
\item SCP03 key objects, \texttt{PUT KEY}, key identifiers and rotated-keyset
      authentication are implemented as regression paths, but diversification,
      production lifecycle governance and official certification vectors still
      remain open work;
\item SCP11a now has an independent establishment vector for P-256, AES-128
      and S16. SCP11b and SCP11c still rely on development-profile campaigns
      for their current cryptographic regression evidence;
\item SCP11a, SCP11b and SCP11c currently have development-profile
      establishment and protected-messaging campaigns using the same
      Amendment~D payload-and-MAC framing as SCP03; those
      campaigns enforce the SCP parameter/tag-\texttt{84} relation, bind
      Security-Domain-owned SIN/SDIN or Card Group ID values, resolve the
      requested ECKA key selector, and exercise multi-command PSO certificate
      transfer;
\item the current development profile accepts a single OCE certificate of at
      most 256 bytes after PSO command-block reassembly.  Intermediate
      certificate chains selected by \texttt{P2.b8} are explicitly profiled
      out with status \texttt{6A86}.
\end{itemize}

The kernel and Rustlet-backed campaigns also exercise exact status words for
malformed CRT and protected-payload encodings, unknown key selectors, certificate failure,
failed host authentication, invalid MAC chains, stale receipts, replay,
profile-dependent lengths, and commands issued outside their required session
state. These negative tests verify that rejected transformations do not advance
MAC chains or encryption counters and that cryptographic secure-messaging
failures require a fresh establishment sequence.

The goal of the current Appendix is therefore precise: document the regression
evidence used by Oxide SE today, explain how SCP03 S8 and S16 differ in
practice, record the SCP11a/SCP11b/SCP11c validation paths, and make clear which part of
the validation is independently reproduced and which part remains future
certification work.
